% Proposition G: the slow geometry (source: PARADIGM.md, Section 6b).
\section{Proposition G: the slow geometry}\label{sec:propG}

This section concerns joint descent, and variable projection (VarPro) in the induced metric.

\paragraph{Setting.}
\begin{itemize}
  \item $f = f(x)$ is time-invariant and $C^2$.
  \item $x_m$ is $T_0$-periodic, with $\omega_{\mathrm{ref}} := 2\pi/T_0$ and occupation measure
    $\rho := \frac{1}{T_0}\int_0^{T_0}\delta_{x_m(t)}\dd t$.
  \item Law (L) acts in the interior, with gains $(g_a, \delta\gamma_c, \delta\gamma_s)$. Put
    $\Gamma := \operatorname{diag}(\gamma_c I_n, \gamma_s)$, acting per kernel, and $G_\xi := \delta\Gamma$. Write the
    geometry as $\xi := (c, s)$.
  \item $\bar L(a,c,s) := \frac12\int\norm{f - \frac1m\sum_i a_i\varphi(\cdot\,;c_i,s_i)}^2\dd\rho$, and
    $M(c,s) := \frac{1}{m^2}\int\varphi\varphi^\T\dd\rho\otimes I_n$ is the Hessian of $\bar L$ in $a$.
  \item $\mathcal{K}$ is the set of geometries $\xi$ on which (H1) holds, and $\mathcal{K}'\Subset\mathcal{K}$ is
    compact. (The set $\mathcal{K}$ is not the functional $\mathcal{G}(\mu)$ of Section~\ref{sec:thmH}.)
  \item $a^*(\xi) := \operatorname{argmin}_a\bar L$, and $L^*(\xi) := \bar L(a^*(\xi),\xi)$ is the profile
    (variable-projection) loss.
  \item Three dimensionless parameters:
    \begin{itemize}
      \item the rate ratio
        $R := g_a m_0\big/\bigl(\delta\cdot\sup_{\mathcal{K}'}\norm{\Gamma^{1/2}\nabla^2L^*\Gamma^{1/2}}\bigr)$;
      \item the slope $S := (\delta/g_a)\cdot\sup_{\mathcal{K}'}\norm{D_\xi a^*\cdot\Gamma^{1/2}}^2$;
      \item the averaging parameter
        $\varepsilon_{\mathrm{avg}} := (m/k_e)\,\lambda_{\max}(G^{1/2}\nabla^2\bar L\,G^{1/2})/\omega_{\mathrm{eff}}$,
        the second ratio in (H2a$'$). The result tables in \code{results/variable_projection} report its inverse, with $\omega_{\mathrm{ref}}$ in place of
        $\omega_{\mathrm{eff}}$.
    \end{itemize}
\end{itemize}

\paragraph{Hypotheses.}
\begin{itemize}[leftmargin=4.2em,labelsep=0.6em]
  \item[(H0)] On the interval considered, the closed-loop and reduced geometries stay in $\mathcal{K}'$. The
    parameters stay in $\operatorname{int}\ThKe$, and $\sup_{\mathcal{K}'}\max_i\norm{a^*_i} < a_{\max} -
    \varepsilon_K$.
  \item[(H1)] $M \succeq m_0 I$ on $\mathcal{K}$.
  \item[(H2a$'$)] \emph{Two scales.} $k_e \gg \omega_{\mathrm{eff}}$ and $\varepsilon_{\mathrm{avg}} \ll 1$, where
    \begin{equation*}
      \omega_{\mathrm{eff}} := \max\Bigl\{\omega_{\mathrm{ref}},\ \frac{\sup_t\norm{\dot x_m}}{s_{\min}},\
      \sup_X\norm{\partial_x(f - \fhat)},\ \frac{\sup_t\norm{\frac{\mathrm{d}}{\mathrm{d}t}r(x_m(t))}}{\sup\norm{r}}
      \Bigr\}.
    \end{equation*}
    The last entry is the bandwidth of the residual along the reference. Also $\sup\norm{r}/k_e \ll s_{\min}$.
  \item[(H2b)] \emph{Rate separation.} $R \gg 1$: the slowest weight rate $g_a m_0 m/k_e$ is much larger than the
    geometry rate $(m\delta/k_e)\norm{\Gamma^{1/2}\nabla^2L^*\Gamma^{1/2}}$.
  \item[(H2c)] \emph{Flat slow manifold.} $S \ll 1$. Since $Da^* = -M^{-1}\partial^2_{a,\xi}\bar L$, a checkable
    sufficient condition is $\delta\cdot\max(\gamma_c,\gamma_s)\cdot\norm{\partial^2_{a,\xi}\bar L}^2 \ll g_a m_0^2$.
\end{itemize}

\medskip\noindent Three remarks on the hypotheses.
\begin{itemize}
  \item \textbf{(H2b) is about $\boldsymbol{\lambda_{\min}}$.} It is a condition on $\lambda_{\min}$ of the feature
    Gram matrix. Since $0 < \varphi \le 1$, $m_0 \le 1/m^2$. So the slowest weight rate is at most $g_a/(k_e m)$, and
    overlapping kernels make (H2b) fail.
  \item \textbf{The ratios are gain choices.} $R$ and $S$ contain neither $k_e$ nor $\omega_{\mathrm{ref}}$. The
    separation between weights and geometry is therefore set by the gains.
  \item \textbf{(H2b) does not imply (H2c).} The product
    \begin{equation*}
      R\cdot S = m_0\cdot\sup\norm{Da^*\Gamma^{1/2}}^2\big/\sup\norm{\Gamma^{1/2}\nabla^2L^*\Gamma^{1/2}}
    \end{equation*}
    does not depend on $g_a$, $\delta$ or $k_e$. It is fixed by the problem and by the shape of $\Gamma$: in the slope configuration
    below, a numerical check gives values from 166 to 1149 over $\gamma_s/\gamma_c\in[10^{-4},10^{4}]$. It is large
    whenever refitting the weights largely compensates a change of geometry.
\end{itemize}

\begin{thm}{Proposition}{G}[formal reduction]
\mbox{}
\begin{enumerate}[label=(\alph*),leftmargin=2.2em]
  \item[\textup{(a)}] \emph{Two scales, under \textup{(H0)} and \textup{(H2a$'$)}.} To leading order in
    $\omega_{\mathrm{eff}}/k_e$ and $\varepsilon_{\mathrm{avg}}$, the period-averaged closed loop is the joint gradient
    flow
    \begin{equation}\label{eq:propG:joint}
      \dot\vartheta_i = -\frac{m}{k_e}\,G\,\nabla_{\vartheta_i}\bar L(\vartheta_1,\dots,\vartheta_m).
    \end{equation}
    Its measure form is $-\frac{1}{k_e}G\nabla(\delta\bar L/\delta\mu)$.
    \begin{itemize}
      \item This is the drift of $\Aavg$ at $\sigma = \tau = 0$, with $\Phi$ in place of $\tilde\Phi$ and $(a,c,s)$
        in place of $(\alpha,\gamma,\kappa)$, by the same averaging computation as Theorem H(i).
      \item It is not literally a special case of H(i): the saturated coordinates there induce a different,
        state-dependent metric.
      \item $\bar L$ is non-increasing along the averaged flow. Along the closed loop, $\bar L$ at period-averaged
        parameters decreases while the averaged dissipation $\frac{m}{k_e}\norm{\nabla\bar L}_G^2$ dominates the
        $O(\omega_{\mathrm{eff}}/k_e)$ corrections of Proposition H.7 and Remark H.7a. Near equilibrium it can
        increase by $O((\omega_{\mathrm{eff}}/k_e)^2)$, as the weights settle at their shifted equilibrium (in a
        numerical check: $+8\cdot10^{-7}$ over 35 periods at $k_e = 20$, and $3.7\cdot10^{-8}$, about 22 times less,
        at $k_e = 80$).
    \end{itemize}

  \item[\textup{(b$'$)}] \emph{Three scales, for the averaged system \textup{(a)}, under \textup{(H0)}, \textup{(H1)},
    \textup{(H2a$'$)} and \textup{(H2b)}.} After a boundary layer, the weights follow the slow manifold:
    \begin{equation}\label{eq:propG:offset}
      \begin{aligned}
        a - a^*(\xi) &= -\frac{k_e}{m g_a}M^{-1}Da^*\cdot\dot\xi + \textup{h.o.t.}\\
        &= \frac{\delta}{g_a}M^{-1}Da^*\Bigl(\Gamma^{-1} + \frac{\delta}{g_a}(Da^*)^\T Da^*\Bigr)^{-1}
        \nabla L^*(\xi) + \textup{h.o.t.}
      \end{aligned}
    \end{equation}
    This reduces to $\frac{\delta}{g_a}M^{-1}Da^*\,\Gamma\,\nabla L^*(\xi)$ when $S \ll 1$. The geometry follows, to
    leading order in $1/R$ and for every $S$,
    \begin{equation}\label{eq:propG:induced}
      \dot\xi = -\frac{m\delta}{k_e}\Bigl(\Gamma^{-1} + \frac{\delta}{g_a}(Da^*)^\T Da^*\Bigr)^{-1}\nabla L^*(\xi).
    \end{equation}
    \begin{itemize}
      \item This is the gradient flow of $(m/k_e)L^*$ in the metric that the $G^{-1}$ metric induces on the graph
        $\{(a^*(\xi),\xi)\}$.
      \item Its metric tensor is $G_\xi^{-1} + g_a^{-1}(Da^*)^\T Da^*$, and it couples all kernels through $Da^*$.
      \item In the closed loop, the weight equilibrium is further shifted by $O(\omega_{\mathrm{eff}}/k_e)$
        (Proposition~H.7). This shift can exceed the offset displayed above: in a numerical check on the
        \code{satisfying} run it is $3$--$6\cdot10^{-3}$, against an offset of $4\cdot10^{-5}$--$3\cdot10^{-4}$.
    \end{itemize}

  \item[\textup{(b)}] \emph{Variable projection, under \textup{(H2c)} in addition.} Let $(\bar c,\bar s)$ solve
    $\bar\xi' = -\frac{m\delta}{k_e}\Gamma\nabla L^*(\bar\xi)$, that is,
    $\bar c_i' = -\frac{m\delta\gamma_c}{k_e}\nabla_{c_i}L^*$ and
    $\bar s_i' = -\frac{m\delta\gamma_s}{k_e}\partial_{s_i}L^*$, with $(\bar c,\bar s)(0) = (c,s)(0)$. Then
    \begin{equation}\label{eq:propG:varpro}
      \sup_{t\le\tau_1k_e/(m\delta)}\norm{(c,s)(t) - (\bar c,\bar s)(t)} \to 0
      \quad\text{as } (\omega_{\mathrm{eff}}/k_e,\ \varepsilon_{\mathrm{avg}},\ S,\ 1/R)\to0,
      \text{ with } \tau_1 \text{ fixed.}
    \end{equation}
    \begin{itemize}
      \item Formally the error is $O(\omega_{\mathrm{eff}}/k_e + \varepsilon_{\mathrm{avg}} + S + 1/R)$.
      \item $L^*$ is non-increasing along $\bar\xi$, and the weights track $a^*$.
      \item This flow is the $S\to0$ limit of (b$'$).
    \end{itemize}
\end{enumerate}
\end{thm}
\status{Formal (a derivation without error control); checked directly, including a counterexample to plain VarPro at
$S = O(1)$ (see the numerics below). Novelty (literature audit): a fixed-$m$, closed-loop analogue of the
two-timescale/VarPro training of~\cite{S1}, \cite{S2}.}

\paragraph{No mean-field form.} For fixed $m$, on configurations of $m$ distinct geometries in $\mathcal{K}$, the flow
(b) is the gradient flow of $L^*/k_e$ in the metric $(1/m)\sum_i\norm{\dot\xi_i}^2_{G_\xi^{-1}}$. Locally, that metric
is the Wasserstein metric restricted to uniform $m$-atom measures.
\begin{itemize}
  \item $L^*$ is the unregularised profile loss. It is the $\lambda\to0$ limit of $\lambda\mathcal{L}^\lambda$
    of~\cite{S2} (time rescaled by $\lambda$). It is not the $\lambda\to0$ functional $\mathcal{L}^0$ of~\cite{S2},
    which is a min-norm interpolation functional whose flow is an ultra-fast diffusion.
  \item As a functional of the geometry marginal, $L^*$ depends only on the support, because the free weights absorb
    the masses. So it is not $W$-continuous, and the statement has no mean-field form. This is consistent with
    $m_0 \le 1/m^2$.
\end{itemize}

\begin{proof}[Derivation (formal)]
Each step is a standard theorem on its own time scale. Their composition on the $O(k_e/(m\delta))$ horizon is formal.
\begin{enumerate}[label=\arabic*.,leftmargin=2em]
  \item \emph{Error channel (exact identity).} For $w := e - r(x,\vartheta)/k_e$, $\dot w = -k_e w - \dot r/k_e$. So
    $\norm{w(t)} \le e^{-k_et}\norm{w(0)} + \sup\norm{\dot r}/k_e^2$, and
    \begin{equation*}
      e = r(x_m)/k_e + O\bigl(\sup\norm{\dot r}/k_e^2 + \sup\norm{\partial_x r}\cdot\sup\norm{r}/k_e^2\bigr).
    \end{equation*}
    This is $O(\omega_{\mathrm{eff}}\cdot\sup\norm{r}/k_e^2)$, with the bandwidth term in $\omega_{\mathrm{eff}}$.
  \item \emph{Period averaging of the joint system} (Khalil~\cite[Thm.~10.4]{khalil2002};
    Sastry--Bodson~\cite[Ch.~4]{sb1989}). This gives (a), on $[0, b/\varepsilon_{\mathrm{avg}}]$.
  \item \emph{Weights against geometry.} Use Tikhonov in slow time (Kokotović--Khalil--O'Reilly~\cite{kko1986}) with
    the uniformly exponentially stable averaged weights on $\mathcal{K}'$ (H1).
    \begin{itemize}
      \item On the slow manifold $a = h(\xi)$,
        $\dot a = Dh\,\dot\xi = -\frac{mg_a}{k_e}M(h - a^*)$, so
        $h - a^* = -\frac{k_e}{mg_a}M^{-1}Dh\,\dot\xi$. This must be evaluated with the (b$'$) velocity, not the
        VarPro velocity.
      \item Since $\bar L$ is quadratic in $a$ and $\partial^2_{\xi a}\bar L = -(Da^*)^\T M$, we get
        $\nabla_\xi\bar L(h,\xi) = \nabla L^* - (Da^*)^\T M(h - a^*) + O((\delta/g_a)^2)$.
      \item Substituting into $\dot\xi = -\frac{m\delta}{k_e}\Gamma\nabla_\xi\bar L(h,\xi)$, with
        $Dh = Da^* + O(1/R)$, gives (b$'$). The resummed form solves the invariance equation exactly once
        $Dh - Da^*$ and the quadratic term are dropped, so the remaining corrections are in $1/R$ (they involve
        $D^2a^*$ and $\partial_\xi M$), not in $S$.
    \end{itemize}
  \item \emph{Envelope theorem.} $\nabla_\xi\bar L(a^*(\xi),\xi) = \nabla L^*(\xi)$. When $S\ll1$, the metric
    correction in (b$'$) drops out, which gives (b).
\end{enumerate}
The composition requires the averaged system to remain valid up to $t\sim\tau_1k_e/(m\delta)$, which is much longer
than $b/\varepsilon_{\mathrm{avg}}$. Step~2 does not certify this. A proof would combine:
\begin{itemize}
  \item a near-identity averaging transform;
  \item uniform exponential stability of the averaged weights on $\mathcal{K}'$;
  \item Tikhonov in slow time with an $O(\varepsilon_{\mathrm{avg}})$ perturbation;
  \item Grönwall in $\xi$ (Sastry--Bodson~\cite[\S4.4]{sb1989}; Sanders--Verhulst--Murdock~\cite[Ch.~5]{svm2007}).
\end{itemize}
\end{proof}

\paragraph{Numerics.} The results are in \code{results/variable_projection} of the accompanying repository;
Table~\ref{tab:propG:numerics} summarises the comparison.
\begin{itemize}
  \item The first two rows use centred differences of period means.
  \item The slope rows use the displacement between the first and last period means. They start on the slow manifold
    of (b$'$), so no weight boundary layer enters the comparison.
  \item (b$'$) was evaluated only in the slope rows.
  \item $k_e/\omega_{\mathrm{eff}}$ is 2.0 in the first two rows with the box's $s_{\min}$, and 2.6--4.8 with the
    actual widths, so (H2a$'$) holds only marginally there. In the \code{satisfying} run the residual bandwidth, about
    7.7, then dominates.
  \item $k_e/\omega_{\mathrm{eff}}$ is about 16 in the slope rows.
\end{itemize}

\begin{table}[htbp]
\centering
\footnotesize
\begin{tabularx}{\textwidth}{@{}>{\raggedright\arraybackslash}p{2.35cm}llYY@{}}
\toprule
configuration & (H2b) ratio & $S$ & G(a): closed loop against the joint flow (a) & G(b) and the induced-metric flow
(b$'$)\\
\midrule
\code{violating} ($m = 4$, overlapping) & 0.10--0.23 & 5.4--25 &
joint cosine 0.9994 / 0.9985, size ratio 1.00 / 0.94 &
(b) fails: cosine $-0.12$ / 0.39, and $L^*$ \emph{increases} from $3.8\cdot10^{-3}$ to $5.9\cdot10^{-3}$\\
\addlinespace
\code{satisfying} ($m = 2$) & 79--87 & 0.006 &
joint cosine 0.985 / 0.962, size 0.96 / 0.94; geometry block 0.999 / 0.998; the weight block fails (see below) &
(b) holds: cosine 0.999, ratio 0.97; $L^*$ decreases from $1.20\cdot10^{-2}$ to $4.11\cdot10^{-3}$\\
\addlinespace
\code{slope}, $\delta = 0.47$ ($m = 1$) & 81--178 & 2.0--2.8 &
cosine 0.9992, size 1.10 &
(b) fails: ratio 0.42, cosine 0.93, and the centre moves the opposite way; (b$'$): cosine 0.9986, ratio 1.12\\
\addlinespace
\code{slope}, $\delta = 0.47$, all gains $\div\,4$ & 81--178 & 2.0--2.8 &
cosine 1.0000, size 1.05 &
(b) fails: ratio 0.40, cosine 0.92; (b$'$): cosine 0.9999, ratio 1.065\\
\addlinespace
\code{slope}, $\delta = 0.047$ & 654--1779 & 0.20--0.32 &
cosine 1.0000, size 1.04 &
(b): ratio 0.85, cosine 0.9986; (b$'$): ratio 1.05, cosine 1.0000\\
\addlinespace
\code{slope}, $\delta = 0.0047$ & 13162--17790 & 0.020--0.023 &
cosine 1.0000, size 1.04 &
(b): ratio 1.02; (b$'$): ratio 1.04; both cosines 1.0000\\
\bottomrule
\end{tabularx}
\caption{Proposition G: the closed loop against the averaged predictions. In the first two rows, an entry
`$p$ / $q$' gives the medians over the first quarter ($p$) and over the last quarter ($q$) of the run.}
\label{tab:propG:numerics}
\end{table}

\begin{itemize}
  \item \textbf{Joint descent.} In all runs, $\bar L$ at the period-mean parameters is non-increasing, sampled once per
    period.
  \item \textbf{(H2b) is not enough.} The \code{slope} rows satisfy (H2b) by a wide margin, yet (b) fails until $S$ is
    small. Meanwhile (b$'$) tracks the closed loop at every $S$. So (H2c) is the hypothesis missing from an earlier
    formulation of Proposition G(b), and (b$'$) is the correct leading-order slow flow. The configuration is a
    counterexample to that earlier formulation, made reproducible in \code{experiments/exp_slow_manifold_slope.py}.
  \item \textbf{Error decomposition} (\code{slope}, $\delta = 0.47$, started on the slow manifold):
    \begin{itemize}
      \item $\norm{(\mathrm{a})}/\norm{(\mathrm{b}')} = 1.015$. This is the intrinsic $O(1/R)$ error of (b$'$).
      \item $\norm{\mathrm{closed}}/\norm{(\mathrm{a})} = 1.10$. It falls to 1.05 when all gains are divided by 4,
        which leaves $R$ and $S$ unchanged and divides $\varepsilon_{\mathrm{avg}}\approx0.08$ by 4. So this part is
        the averaging error.
      \item The remaining ${\approx}\,4.4\%$ is the $O(1/k_e)$ error-channel factor at $k_e = 20$. It is the same at
        $\delta = 0.047$ and $0.0047$.
    \end{itemize}
  \item \textbf{The size ratio in the \code{satisfying} run.} The centred ratio of 0.97 is the error-channel factor
    times a $+1.4\%$ averaging correction from the weights' oscillation within each period, the analogue of Remark
    H.7a.
    \begin{itemize}
      \item The frozen-parameter ratio is 0.951--0.956 along the run, and 0.996 at $k_e = 80$.
      \item The fastest weight rate over $\omega_{\mathrm{ref}}$ is about 0.14.
      \item In a further numerical check, with the weights running at frozen geometry the ratio is 0.969--0.970;
        with $g_a/4$ it is 0.956--0.960.
      \item The slope correction is below 0.01\% here, because $\nabla L^*$ is nearly orthogonal to the range of
        $(Da^*)^\T$. An earlier estimate of 0.94--0.95 contained a 2.5\% forward-difference bias.
    \end{itemize}
  \item \textbf{The weight block.} In the \code{satisfying} run the weights sit at $a^* + O(1/k_e)$, within 0.1\% of
    $a^*$.
    \begin{itemize}
      \item Their actual drift is tiny: 0.001--0.008 of the geometry speed in the $G^{-1}$ norm.
      \item The prediction $-(m/k_e)g_a\nabla_a\bar L$ at the period-mean weights is 0.17--0.30 of the geometry speed,
        40--200 times the actual drift. It is the restoring velocity of the $O(1/k_e)$ offset, which (a) does not
        contain.
      \item So the weight block fails (cosine $-0.21$ / 0.50, carrying 92--95\% of the residual), while the geometry
        block matches.
    \end{itemize}
\end{itemize}
\checkedby{\code{experiments/exp_variable_projection.py}, which writes \code{table_violating.md} and
\code{table_satisfying.md}; and \code{experiments/exp_slow_manifold_slope.py}, which writes \code{table_slope.md}
(all in \code{results/variable_projection}).}

\paragraph{Status.}
\begin{itemize}
  \item \textbf{What the loop supplies.} The error-channel and averaging scales, and the data distribution (the
    occupation measure of the reference), come from the closed loop.
  \item \textbf{What the designer supplies.} The separation between weights and geometry is a gain choice, as
    in~\cite{S1}, \cite{S2}:
    $\delta \ll g_a\cdot\min\bigl(m_0/\norm{\Gamma^{1/2}\nabla^2L^*\Gamma^{1/2}},\
    1/\norm{Da^*\Gamma^{1/2}}^2\bigr)$. Equivalently, $\norm{G_\xi^{1/2}\nabla^2L^*G_\xi^{1/2}} \ll g_a m_0$ and
    $\norm{Da^*G_\xi^{1/2}}^2 \ll g_a$.
  \item \textbf{What is new here.}
    \begin{itemize}
      \item The learner's data distribution is its own excitation, and at finite $k_e$ it moves with the parameters
        (Proposition H.7(c)).
      \item The slow flow is variable projection in the \emph{induced} metric, valid for every $S$. The slope $S$
        measures the departure from plain VarPro.
    \end{itemize}
\end{itemize}
